% =====================================================================
% แนวทางการเขียนเล่มคู่มือการใช้งาน (สำหรับนักศึกษา)
% หน้านี้แสดงเฉพาะเมื่อใช้ตัวเลือก guides ใน \documentclass[guides]{cskkumanual}
% ลบตัวเลือก guides ออกก่อนส่งเล่ม หน้านี้และกล่อง "แนวทางการเขียน" ทุกกล่องจะหายไปเอง
% =====================================================================
เล่มคู่มือการใช้งานเป็นเอกสารที่ส่งพร้อมระบบ เพื่อให้ผู้อื่นติดตั้ง ใช้งาน ดูแล และพัฒนาระบบต่อได้
โดยไม่ต้องถามผู้จัดทำ ต่างจากเล่มรายงานโครงงานซึ่งอธิบายว่าโครงงานทำอะไร ทำไม และได้ผลอย่างไร
ตัวอย่างในเล่มนี้เป็นคู่มือของระบบ CS KKU Project Manager
ให้เขียนตามโครงเดียวกันแล้วแทนเนื้อหาด้วยระบบของตนเอง
และอ่านกล่อง ``แนวทางการเขียน'' ที่ต้นแต่ละบทประกอบ

\section*{ผู้อ่านของคู่มือ}

คู่มือหนึ่งเล่มมีผู้อ่าน 3 กลุ่ม แต่ละกลุ่มต้องการข้อมูลต่างกัน จึงแยกบทตามกลุ่มผู้อ่าน ดังนี้
% ตารางในหน้านี้ไม่มีชื่อตาราง เพื่อไม่ให้เลขตารางของคู่มือเลื่อน

\begin{center}
\begin{tabular}{|L{3.3cm}|L{\dimexpr\linewidth-6.3cm-6\tabcolsep-4\arrayrulewidth\relax}|L{3cm}|}
\hline
\textbf{ผู้อ่าน} & \textbf{ต้องการรู้} & \textbf{บทที่เกี่ยวข้อง} \\ \hline
ผู้ใช้งานทั่วไป & วิธีใช้ระบบทำงานของตนให้เสร็จ ตามบทบาทของตน & 1, 3--5, 8 \\ \hline
ผู้ดูแลระบบ & ความต้องการของระบบ การติดตั้ง ตั้งค่า สำรองข้อมูล และปรับรุ่น & 2, 6, 8 \\ \hline
นักพัฒนา & สถาปัตยกรรม โครงสร้างโค้ด ฐานข้อมูล API และวิธีพัฒนาต่อ & 2, 7, ภาคผนวก \\ \hline
\end{tabular}
\end{center}

\section*{โครงสร้างเล่มที่แนะนำ}

\addtocounter{table}{-1}% longtable ไม่มีชื่อตาราง: ไม่นับเลขตาราง
\begin{longtable}{|L{4.4cm}|L{\dimexpr\linewidth-8.4cm-6\tabcolsep-4\arrayrulewidth\relax}|L{4cm}|}
\hline
\textbf{ส่วน/บท} & \textbf{เนื้อหาหลัก} & \textbf{ต้องมีเมื่อ} \\ \hline
\endhead
คำนำ ประวัติการแก้ไขเอกสาร & ระบบอะไร รุ่นใด ใครควรอ่าน และการแก้ไขคู่มือแต่ละครั้ง & ทุกโครงงาน \\ \hline
1 ภาพรวมของระบบ & ระบบทำอะไร ผู้ใช้และสิทธิ์ วิธีอ่านคู่มือ สัญลักษณ์ในคู่มือ & ทุกโครงงาน \\ \hline
2 ความต้องการของระบบและการติดตั้ง & ฮาร์ดแวร์ ซอฟต์แวร์ ขั้นตอนติดตั้ง การตั้งค่า การตรวจสอบหลังติดตั้ง & ทุกโครงงาน \\ \hline
3 การเริ่มต้นใช้งาน & สมัครสมาชิก เข้าสู่ระบบ ส่วนประกอบหน้าจอหลัก & มีส่วนติดต่อกับผู้ใช้ \\ \hline
4--5 การใช้งานตามบทบาท (คู่มือผู้ใช้งานทั่วไป) & หนึ่งบทต่อหนึ่งบทบาท หนึ่งหัวข้อต่อหนึ่งงาน & มีผู้ใช้หลายบทบาท (บทบาทเดียวรวมเป็นบทเดียว) \\ \hline
6 การดูแลระบบ & งานประจำ ผู้ใช้และสิทธิ์ สำรอง/กู้คืนข้อมูล ปรับรุ่น ตรวจสถานะ ความปลอดภัย & มีเครื่องแม่ข่ายหรือฐานข้อมูล \\ \hline
7 คู่มือสำหรับนักพัฒนา & สถาปัตยกรรม โครงสร้างโค้ด สภาพแวดล้อมพัฒนา ฐานข้อมูล API แนวทางพัฒนาต่อ การนำขึ้นใช้งาน & ทุกโครงงาน \\ \hline
8 การแก้ไขปัญหาเบื้องต้น & อาการ--สาเหตุ--วิธีแก้ คำถามที่พบบ่อย ช่องทางติดต่อ & ทุกโครงงาน \\ \hline
ภาคผนวก & อภิธานศัพท์ รายการ API พจนานุกรมข้อมูล & ตามความเหมาะสม \\ \hline
\end{longtable}

\section*{ปรับตามประเภทโครงงาน}

โครงเล่มเหมือนกันทุกประเภท แต่เนื้อหาในบทต่างกัน ดังนี้

\addtocounter{table}{-1}%
\begin{longtable}{|L{2.6cm}|L{\dimexpr0.5\linewidth-1.3cm-3\tabcolsep-2\arrayrulewidth\relax}|L{\dimexpr0.5\linewidth-1.3cm-3\tabcolsep-2\arrayrulewidth\relax}|}
\hline
\textbf{ประเภท} & \textbf{ความต้องการของระบบและการติดตั้ง} & \textbf{การใช้งานและคู่มือนักพัฒนา} \\ \hline
\endhead
เว็บแอปพลิเคชัน & เบราว์เซอร์ที่รองรับ สเปกเครื่องแม่ข่าย ฐานข้อมูล การติดตั้งด้วย Docker หรือสคริปต์ การตั้งค่าเว็บเซิร์ฟเวอร์
  & ภาพหน้าจอตามบทบาท \newline API ฐานข้อมูล และการนำขึ้นใช้งาน \\ \hline
แอปพลิเคชันบนมือถือ & รุ่น Android/iOS ที่รองรับ พื้นที่ว่าง สิทธิ์ที่แอปขอ วิธีติดตั้งไฟล์ APK หรือจากสโตร์ และส่วนหลังที่แอปเชื่อมต่อ
  & ภาพหน้าจอมือถือ ท่าทางสัมผัส การแจ้งเตือน \newline การ build และ signing การตั้งค่า endpoint \\ \hline
โปรแกรมบนคอมพิวเตอร์ & ระบบปฏิบัติการ runtime ที่ต้องมี วิธีติดตั้งจากไฟล์ติดตั้ง และตำแหน่งไฟล์ตั้งค่า
  & เมนูและหน้าต่างของโปรแกรม \newline การ build ไฟล์ติดตั้ง \\ \hline
ระบบฝังตัว/IoT & รายการอุปกรณ์พร้อมรุ่น แผนผังการต่อวงจร แหล่งจ่ายไฟ การอัปโหลดเฟิร์มแวร์ และการตั้งค่าเครือข่าย
  & การเปิดใช้งาน ความหมายของไฟสถานะ ข้อควรระวังด้านไฟฟ้า \newline การกำหนดขา (pin) โปรโตคอลสื่อสาร และโค้ดเฟิร์มแวร์ \\ \hline
งานวิจัย/โมเดล AI & รุ่น Python ไลบรารี GPU/CUDA ตำแหน่งดาวน์โหลดโมเดลและชุดข้อมูล
  & วิธีเรียกใช้โมเดล รูปแบบข้อมูลนำเข้าและผลลัพธ์ \newline การเตรียมข้อมูล การฝึกโมเดลใหม่ และการวัดผล \\ \hline
\end{longtable}

\section*{หลักการเขียน}

\begin{enumerate}
  \item เขียนจากมุมของผู้อ่าน ใช้ชื่อปุ่ม เมนู และข้อความให้ตรงกับที่ปรากฏบนหน้าจอ
  \item หนึ่งหัวข้อคือหนึ่งงาน ตั้งชื่อหัวข้อขึ้นต้นด้วย ``การ...'' เขียนเป็นขั้นตอน 1-2-3
        และบอกผลลัพธ์ที่ควรเห็นหลังขั้นตอนสุดท้าย
  \item ภาพหน้าจอครอปเฉพาะส่วนที่อธิบายให้ตัวอักษรอ่านออก ใส่หมายเลขกำกับและอธิบายทุกหมายเลข
        ใช้ข้อมูลทดสอบในภาพ ไม่ใช้ข้อมูลจริงของบุคคล
  \item คำสั่งและค่าตั้งค่าต้องคัดลอกไปใช้ได้ทันที และระบุรุ่นของซอฟต์แวร์ทุกตัว
  \item ห้ามใส่รหัสผ่านจริง คีย์ API หรือข้อมูลลับ ให้ใช้ค่าตัวอย่างและบอกว่าต้องเปลี่ยน
  \item ไม่คัดลอกความเป็นมา ทฤษฎี หรือแผนภาพทั้งหมดจากเล่มรายงานมาซ้ำ ให้อ้างถึงเล่มรายงานแทน
  \item ไม่ใส่ภาพหน้าจอการติดตั้งโปรแกรมทั่วไป (เช่น VS Code, XAMPP) ทีละหน้า ให้ระบุชื่อ รุ่น และแหล่งดาวน์โหลด
  \item ทดสอบคู่มือก่อนส่ง: ให้ผู้ที่ไม่เคยใช้ระบบติดตั้งตามบทที่ 2 บนเครื่องใหม่
        และทำงานตามคู่มือผู้ใช้โดยไม่มีผู้จัดทำช่วย บันทึกจุดที่ติดแล้วแก้คู่มือ
  \item แก้คู่มือทุกครั้งที่ระบบเปลี่ยน และบันทึกในประวัติการแก้ไขเอกสาร
\end{enumerate}

\section*{ตรวจก่อนส่งเล่ม}

\begin{itemize}[label=$\square$]
  \item ลบตัวเลือก \cmd{guides} ใน \filename{main.tex} และลบเนื้อหาตัวอย่างของระบบ CS KKU Project Manager ออกหมด
  \item ชื่อปุ่ม เมนู ภาพหน้าจอ และรุ่นของระบบในคู่มือตรงกับระบบรุ่นที่ส่ง
  \item ทุกคำสั่งในบทที่ 2 และบทที่ 7 ทดลองบนเครื่องใหม่แล้ว
  \item ภาพและตารางทุกรายการมีชื่อ และอ้างถึงในเนื้อหา
  \item ไม่มีรหัสผ่านจริง คีย์ หรือข้อมูลส่วนบุคคลในภาพและตัวอย่าง
\end{itemize}
